% Blueprint for the extremal half of Erdős problem 1007.
%
% Every node below is grounded in Chaffee and Noble, "Dimension 4 and dimension 5 graphs with
% minimum edge set", Australas. J. Combin. 64(2) (2016), 327-333, read in full from
% https://ajc.maths.uq.edu.au/pdf/64/ajc_v64_p327.pdf. Citations are to that paper's own numbering.
% Where a node has no counterpart in the paper it is marked GAP and says what the paper omits.
%
% `\lean{...}` names a declaration that must exist; `leanblueprint checkdecls` verifies it.
% Never write a `\lean{}` for a declaration that has not been built.

\chapter{A dimension four graph with nine edges is \texorpdfstring{$K_{3,3}$}{K33}}

\section{Definitions}

\begin{definition}[Unit-distance representation]
  \label{def:representable}
  \lean{Erdos1007.Standalone.Mathlib.InlineErdos1007.UnitDistanceEmbeddable}
  \leanok
  A graph $G$ is \emph{representable in $\mathbb{R}^n$} when there is an injective
  $f : V(G) \to \mathbb{R}^n$ with $\lVert f(u) - f(v)\rVert = 1$ for every edge $uv$.
  Non-adjacent vertices are unconstrained: they may also lie at distance one.
\end{definition}

\begin{definition}[Dimension]
  \label{def:dimension}
  \lean{Erdos1007.Standalone.Mathlib.InlineErdos1007.HasDimension}
  \leanok
  \uses{def:representable}
  $\dim(G) = n$ when $n$ is \emph{least} such that $G$ is representable in $\mathbb{R}^n$.
\end{definition}

\section{Imported results}

Chaffee--Noble import these from Erd\H{o}s, Harary and Tutte, \emph{On the dimension of a graph},
Mathematika \textbf{12} (1965), 118--122, and use them without proof. Only the three below are
needed here. They are obligations of this development, not givens.

\begin{lemma}[Monotonicity; EHT via Chaffee--Noble Lemma 4]
  \label{lem:monotone}
  If $H$ is a subgraph of $G$ then $\dim(H) \le \dim(G)$.
\end{lemma}

\begin{lemma}[$\dim(K_n - e) = n - 2$; EHT via Chaffee--Noble Lemma 2]
  \label{lem:kn-minus-e}
  Only the instance $\dim(K_5 - e) = 3$ is used, and only the upper bound $\le 3$.
\end{lemma}

\begin{lemma}[$\dim(K_{3,3}) = 4$; EHT via Chaffee--Noble Lemma 3]
  \label{lem:k33-dim}
  \uses{def:dimension}
  Both halves are needed and neither is quoted in Chaffee--Noble.

  \emph{Upper bound.} With $r = 1/\sqrt{2}$, the parts
  $\{(r,0,0,0), (-r,0,0,0), (0,r,0,0)\}$ and $\{(0,0,r,0), (0,0,-r,0), (0,0,0,r)\}$ give six
  distinct points whose nine cross-distances are all one.

  \emph{Lower bound.} Suppose $K_{3,3}$ were representable in $\mathbb{R}^3$ with parts $A, B$.
  Each $b \in B$ lies on all three unit spheres centred at the points of $A$. If those centres are
  non-collinear the three spheres meet in at most two points, so three distinct elements of $B$
  cannot all lie there. If they are collinear and distinct, a common point would lie on two
  distinct parallel bisector planes, which is impossible. Either way $\lvert B\rvert \le 2$.
\end{lemma}

\section{Minimum degree}

\begin{lemma}[Few edges force dimension at most three]
  \label{lem:eight-edges}
  \uses{def:dimension, lem:monotone, lem:kn-minus-e}
  Every graph with at most eight edges has dimension at most three.

  This is the content of Chaffee--Noble Theorem 6. Take $G$ with $\dim(G) > 3$ and $|E(G)|$
  minimum, and suppose $|E(G)| \le 8$. Minimality excludes a vertex of degree one. It also
  excludes a vertex $u$ of degree two: writing $uv_1, uv_2$ for its edges, set
  $V(G') = V(G)\setminus\{u\}$ and $E(G') = (E(G)\setminus\{uv_1,uv_2\})\cup\{v_1v_2\}$. Then
  $|E(G')| < |E(G)|$, so $G'$ embeds in $\mathbb{R}^3$; there $v_1$ and $v_2$ are at distance one,
  so the unit spheres about them meet in a circle, and any point of that circle serves as $u$.
  That embeds a supergraph of $G$, so $\dim(G) \le 3$ by Lemma~\ref{lem:monotone}.

  Adding edges until $|E(H)| = 8$ with every degree at least three, the degree sum is sixteen, so
  the degree sequence is $(4,3,3,3,3)$; hence $H \subseteq K_5 - e$ and
  $\dim(H) \le 3$ by Lemmas~\ref{lem:kn-minus-e} and~\ref{lem:monotone}.
\end{lemma}

\begin{lemma}[Minimum degree three]
  \label{lem:min-degree}
  \uses{lem:eight-edges, lem:monotone}
  If $\dim(G) = 4$ and $|E(G)| = 9$ then every vertex of $G$ has degree at least three.

  Chaffee--Noble discharge this in Theorem 7 with the phrase ``by the arguments in the preceding
  theorem''. Deleting a vertex of degree at most two leaves at most seven edges, so
  Lemma~\ref{lem:eight-edges} applies and the re-placement argument of
  Lemma~\ref{lem:eight-edges} puts $G$ in $\mathbb{R}^3$.
\end{lemma}

\section{Bounding the vertex count}

\begin{lemma}[At least five vertices]
  \label{lem:at-least-five}
  A graph with nine edges has at least five vertices, since $K_4$ has only six.
\end{lemma}

\begin{lemma}[At most six vertices]
  \label{lem:at-most-six}
  \uses{lem:min-degree}
  If every degree is at least three and $|E(G)| = 9$ then $|V(G)| \le 6$, because
  $3\lvert V(G)\rvert \le \sum_{v} \deg v = 2\lvert E(G)\rvert = 18$.

  \textbf{GAP.} This step does not appear in Chaffee--Noble. The paper establishes the lower bound
  on $|V|$ and then treats five and six vertices as the only cases, without saying why larger
  vertex counts are excluded. The handshake bound above supplies it. Any development that claims
  to follow the paper step by step must insert this.
\end{lemma}

\section{The five-vertex case}

\begin{lemma}[Five vertices is impossible]
  \label{lem:five-vertices}
  \uses{lem:min-degree, lem:kn-minus-e, lem:monotone}
  There is no $G$ with $\dim(G) = 4$, $|E(G)| = 9$ and $|V(G)| = 5$.

  Degrees lie between three and four, and the degree sum is eighteen, forcing $(4,4,4,3,3)$.
  Chaffee--Noble assert that this sequence ``corresponds solely to the graph $K_5 - e$''; that
  uniqueness is stated without argument and must be proved here. Then
  $\dim(K_5 - e) = 3$ contradicts $\dim(G) = 4$.
\end{lemma}

\section{The six-vertex case}

\begin{lemma}[Two-regular graphs on six vertices]
  \label{lem:two-regular-six}
  A two-regular graph on six vertices is isomorphic to $C_6$ or to $K_3 \sqcup K_3$.

  Chaffee--Noble state this as evident. It is the combinatorial core of the formalization and has
  no ready route in Mathlib.
\end{lemma}

\begin{lemma}[The prism embeds in the plane]
  \label{lem:prism-planar}
  \uses{def:representable}
  The complement of $C_6$ is representable in $\mathbb{R}^2$, hence has dimension at most two.

  Chaffee--Noble give the placement
  $(0,0),\ (1,0),\ (0,1),\ (1,1),\ (\tfrac12, \tfrac{\sqrt3}{2}),\ (\tfrac12, 1 + \tfrac{\sqrt3}{2})$.
  Labelling these $A,B,C,D,E,F$, the nine unit distances are $AB, CD, AC, BD, AE, BE, CF, DF, EF$:
  triangles $ABE$ and $CDF$ joined by the matching $AC, BD, EF$, which is the triangular prism.
  Verified by direct computation against the paper.
\end{lemma}

\begin{lemma}[Six vertices forces $K_{3,3}$]
  \label{lem:six-vertices}
  \uses{lem:min-degree, lem:two-regular-six, lem:prism-planar}
  If $\dim(G) = 4$, $|E(G)| = 9$ and $|V(G)| = 6$ then $G \cong K_{3,3}$.

  Minimum degree three and degree sum eighteen force three-regularity, so the complement is
  two-regular. By Lemma~\ref{lem:two-regular-six} the complement is $C_6$ or $K_3 \sqcup K_3$. The
  first makes $G$ the prism, of dimension at most two by Lemma~\ref{lem:prism-planar},
  contradicting $\dim(G) = 4$. The second makes $G$ exactly $K_{3,3}$.
\end{lemma}

\section{The target}

\begin{theorem}[Chaffee--Noble, Theorem 7]
  \label{thm:target}
  \lean{Erdos1007.Standalone.Mathlib.InlineErdos1007.DimensionFourExtremal}
  \uses{lem:min-degree, lem:at-least-five, lem:at-most-six, lem:five-vertices, lem:six-vertices}
  A graph of dimension four with nine edges and no isolated vertex is isomorphic to $K_{3,3}$.
\end{theorem}

\begin{proof}
  \uses{lem:min-degree, lem:at-least-five, lem:at-most-six, lem:five-vertices, lem:six-vertices}
  Lemma~\ref{lem:min-degree} gives minimum degree three, so
  Lemmas~\ref{lem:at-least-five} and~\ref{lem:at-most-six} leave five or six vertices.
  Lemma~\ref{lem:five-vertices} excludes five and Lemma~\ref{lem:six-vertices} settles six.
\end{proof}

\begin{theorem}[The statement is not vacuous]
  \label{thm:witness}
  \lean{Erdos1007.Standalone.Mathlib.InlineErdos1007.DimensionFourExtremal.witness}
  \uses{lem:k33-dim}
  Some graph has dimension four, nine edges and no isolated vertex, namely $K_{3,3}$ by
  Lemma~\ref{lem:k33-dim}.
\end{theorem}

\section{Separating examples}

\begin{theorem}[Representation does not constrain non-edges]
  \label{thm:sep-representable}
  \lean{Erdos1007.Standalone.Mathlib.InlineErdos1007.UnitDistanceEmbeddable.separating}
  \uses{def:representable}
  Some placement satisfying Definition~\ref{def:representable} puts a non-edge at distance one:
  send an edge to $0, 1$ and a third vertex to $2$ in $\mathbb{R}^1$.
\end{theorem}

\begin{theorem}[Leastness is not membership]
  \label{thm:sep-dimension}
  \lean{Erdos1007.Standalone.Mathlib.InlineErdos1007.HasDimension.separating}
  \uses{def:dimension}
  Some graph with an edge is representable in $\mathbb{R}^4$ without having dimension four, namely
  $K_2$, of dimension one. The edge is required: the empty graph would satisfy this degenerately.
\end{theorem}
